\documentclass[11pt, reqno]{article}

\usepackage{amsmath, amsthm, amssymb}
\usepackage{enumitem}
\usepackage{tcolorbox}
\usepackage{hyperref}
\usepackage{tikz}
\usepackage{tikz-cd}
\usepackage{pgfplots}
\pgfplotsset{compat=1.18}
\usetikzlibrary{arrows.meta}
\usepackage{mathrsfs}
\usepackage{fancyhdr}
\usepackage[bottom=0.75in, top=1in, left=0.5in, right=0.5in]{geometry}
\usepackage{array}   % for \newcolumntype macro
\newcolumntype{L}{>{$}l<{$}}

\theoremstyle{plain}
\newtheorem*{theorem}{Theorem}
\newtheorem*{proposition}{Proposition}
\newtheorem{exercise}{Exercise}
\newtheorem*{lemma}{Lemma}
\newtheorem*{corollary}{Corollary}

\theoremstyle{definition}
\newtheorem*{definition}{Definition}
\newtheorem*{example}{Example}

\theoremstyle{remark}
\newtheorem*{remark}{Remark}

\renewcommand{\phi}{\varphi}
\renewcommand{\epsilon}{\varepsilon}
\renewcommand{\emptyset}{\varnothing}

\newcommand{\RR}{\mathbb{R}}
\newcommand{\ZZ}{\mathbb{Z}}
\newcommand{\NN}{\mathbb{N}}
\newcommand{\CC}{\mathbb{C}}
\newcommand{\QQ}{\mathbb{Q}}

\DeclareMathOperator{\ima}{\text{im}}

\begin{document}

\topmargin=-40pt
\rhead{Henry Woodburn}
\lhead{Math 658}
\renewcommand{\headrulewidth}{1pt}
\renewcommand{\headsep}{20pt}
\thispagestyle{fancy}

{\Huge \bfseries \noindent Homework 1}

\begin{enumerate}
    \item[2.1.1.a] Let $\mu$ be a positive inner regular locally finite measure. Suppose there exists
    a constant $D$ such that for all $x \in \RR^n$ and $r > 0$ we have
    \[
        \mu(3B(x,r)) \leq D \mu(B(x,r)).
    \]
    Define the weak $L^{1,\infty}$ norm $\|\cdot\|_{L^{1, \infty}}$ by
    \[
        \|f\|_{L^{1,\infty}} = \sup_{\lambda > 0}\lambda\mu({|f| > \lambda}).
    \]
    We will show the uncentered maximal function $M$ is continuous $L^1(\RR^n, \mu) \to L^{1,\infty}(\RR^n,\mu)$.
    Define $E_\alpha = \{x \in \RR^n: Mf(x) > \alpha\}$. We need to show
    \[
        \|Mf\|_{L^{1,\infty}} \leq C\|f\|_{L^1},
    \]
    and it will suffice to show for some $\alpha > 0$ that
    \[
        \alpha\mu(E_\alpha) \leq C\|f\|_{L^1}.
    \]

    First, I claim $E_\alpha$ is open. For $x \in E_\alpha$, there must be 
    a ball $B_x$ containing $x$ such that the average value of $f$ on $B_x$ is greater than $\alpha$. Then this ball 
    is in fact contained in $E_\alpha$, since for any point $y \in B_x$ we can choose a ball $B_y \subset B_x$ containing
    $y$, and this proves $Mf(y) > \alpha$. 

    Now use inner regularity. Choose any compact $K \subset E_\alpha$ and cover $K$ in balls $\{B_x\}_{x \in K}$ 
    such the average value of $f$ on every $B_x$ is greater than $\alpha$. Choose a finite subcover $\{B_i\}_1^N$
    and choose a further subcover $\{B_{i_k}\}_1^M$ such that $\cup B_i \subset \cup 3 B_{i_k}$. Then
    \[
        \mu(K) \leq \mu\left(\bigcup_{k = 1}^M 3 B_{i_k}\right) \leq D \sum \mu(B_{i_k}) \leq \frac{D}{\alpha} \sum \int_{B_{i_k}} |f(y)|dy \leq \frac{D}{\alpha}\|f\|_{L^1}.
    \]

    \item[2.1.2.a] Let $\mathscr{F}$ be a finite family of open intervals in $\RR$. We prove there exist two subfamilies $\mathscr{F}_1$ and $\mathscr{F}_2$
    of $\mathscr{F}$, each family pairwise disjoint, such that the union of the sets in both families is equal to the union of $\mathscr{F}$.

    Without loss of generality, we can reduce to the case when $\cup \mathscr{F}$ is connected, and we can remove any intervals
    strictly contained in larger ones. If the union is not disjoint, we can just solve the problem on disjoint pieces and combine
    the individual families in any way. 

    We proceed by induction. The case of $n = 1$ is trivial. Now assume we have proven that for any family of $k-1$ elements,
    we can find two disjoint subfamilies as above. 

    Let $\mathscr{F}$ be a connected family of $k$ intervals. Remove the leftmost set $U_0$, which must exist since $\mathscr{F}$ is finite, and we 
    have removed any sets contained in others. Then from $\mathscr{F}\setminus\{U_0\}$, choose subfamilies $\mathscr{G}_i$, $i = 1,2$
    which satisfy the desired property.
\end{enumerate}

\end{document}